\chapter{人形动量全身控制：CMP、角动量任务与双支撑过渡}
\label{ch:humanoid-momentum-wbc}

% ════════════════════════════════════════════════════════════════
%  人形（双足平足/窄脚掌）动量全身控制：在 ZMP/DCM/捕获点（\cref{ch:legged-biped}）
%  与 CMM/角动量守恒（\cref{ch:unified-multimodal-mpc}）、WBC-QP（\cref{ch:force-wbc}）
%  之上，组装"角动量任务"，并讲清 CMP、飞轮等效捕获点、双支撑过渡三件人形特有的事。
%  记号归一：向量粗体小写、矩阵粗体大写；h_G=A_G v（CMM 即 A_G）；M v̇+h=Sᵀtau+ΣJ_cᵀlambda。
% ════════════════════════════════════════════════════════════════

\begin{practice}[前置自测]
开始之前，先试答下列问题。它们都来自“把点足/四足那套落脚平衡，升级到一台\emph{有脚掌、会摆臂}的人形”时真实会卡住的地方；若已能流畅作答，可只速读本章表格与速查卡；若卡住，正是本章要为你解决的。
\begin{enumerate}
  \item \cref{ch:legged-biped} 里 ZMP 必须落在支撑域内，否则物理不可实现。可人在窄独木桥上被推一下时，会猛地甩动手臂稳住——这一甩，是把“地面反力的等效作用点”移到了脚掌\emph{外}吗？如果是，它和那个“不能出脚”的 ZMP 是同一个点吗？
  \item 一个直立的人受到侧推，他有三档反应：先用脚踝微调、不够就甩臂扭腰、再不够就迈一步。这三档在“地面反力等效作用线落在哪”这件事上，几何区别是什么？能否用一个量把“我现在用的是哪一档”读出来？
  \item \cref{ch:unified-multimodal-mpc} 给了角动量守恒 $\dot{\mathbf k}_G=\mathbf 0$（重力不产生绕质心力矩）。可双足平衡里我们又说“摆臂能产生 $\dot{\mathbf k}_G\ne 0$ 来救场”。这两句话矛盾吗？臂的角动量是从哪来的、又去了哪？
  \item 四足 WBC（\cref{ch:quad_wbc}）的任务栈里通常\emph{没有}一条“角动量任务”，照样跑得很好。人形为什么\emph{必须}补上这一条？这条任务的 Jacobian 是什么矩阵的哪一部分？
  \item 双足都有一段两脚同时着地的“双支撑相”。若规划让 ZMP 参考在相位切换瞬间从前脚“啪”地跳到后脚，会在关节力矩上看到什么？为什么一段平滑的 ZMP 转移就能消掉它？
  \item 你用单支撑的解析 DCM 公式去算双支撑期间的落脚，发现总有几厘米的系统偏差。这个偏差的物理根源是什么？它藏在哪个参数里？
\end{enumerate}
\end{practice}

\section*{本章目标}
学完本章，你应当能够：
\begin{enumerate}
  \item \textbf{定义}质心力矩枢轴点 CMP，写出它与 ZMP 的关系 $\mathbf r_{\mathrm{CMP}}=\mathbf r_{\mathrm{ZMP}}+\frac{1}{f_z}[\tau_{G,y};-\tau_{G,x}]$，并\emph{解释}为何 CMP 可以跑到脚掌外、ZMP 不能（\cref{sec:hmw-cmp}）；
  \item \textbf{把}踝、髋/臂、迈步三级恢复策略\emph{几何化}——用“CMP 是否离开 ZMP、捕获点是否出脚”把“机器人正在用哪一档”读成一张平衡仪表盘（\cref{sec:hmw-cmp}）；
  \item \textbf{建立}飞轮模型，理解角动量如何把\emph{等效}捕获点拉近 $\boldsymbol\xi_{eq}=\boldsymbol\xi-\Delta\mathbf k_G/(m\,\omega_0 z_c)$，从而“用角动量买时间、暂时撑大不迈步的可恢复区”（\cref{sec:hmw-flywheel}）；
  \item \textbf{从}质心动量矩阵 CMM $\mathbf h_G=\mathbf A_G\mathbf v$（\cref{ch:unified-multimodal-mpc}）出发，把\emph{角动量任务}推成 WBC-QP 里的一条标准任务 $\mathbf A_G^{\mathrm{ang}}\dot{\mathbf v}=\dot{\mathbf k}_G^{d}-\dot{\mathbf A}_G^{\mathrm{ang}}\mathbf v$（\cref{sec:hmw-momentum-wbc}）；
  \item \textbf{装配}人形 TSID/HQP 的任务栈——支撑接触（硬）、CoM、角动量、摆脚、姿态——并\emph{说清}它比四足 WBIC（\cref{ch:quad_wbc}）多出的正是那条角动量任务（IHMC/TALOS 风格的动量优先架构）（\cref{sec:hmw-momentum-wbc}）；
  \item \textbf{把}双支撑相的 ZMP 参考写成连续转移 $\mathbf p_{\mathrm{zmp}}=(1-s)\mathbf p_{\mathrm{foot}}^{\mathrm{prev}}+s\,\mathbf p_{\mathrm{foot}}^{\mathrm{next}}$、用三/五次多项式保力矩连续，并\emph{说明}单支撑解析 DCM 在双支撑期 $2$–$5\,\mathrm{cm}$ 误差的等效 $\omega_0$ 根源（\cref{sec:hmw-ds}）。
\end{enumerate}

\subsection*{本章知识导航}
本章站在一个具体落差上：\cref{ch:legged-biped} 把\emph{点足}的落脚平衡讲透了，但它有意忽略了两件人形特有的事——\textbf{脚掌有面积、上体会摆臂}。脚掌有面积，意味着踝部能产生力矩、压心（CoP）可在脚掌内移动，于是有了“不迈步也能调”的余地；上体会摆臂，意味着可以主动动用\emph{角动量}去救平衡。这两件事合起来，把“平衡只有迈步一条路”升级成“踝→髋/臂→迈步”的三级策略，并要求控制器在全身 QP 里显式管住角动量。本章的灵魂就是这条升级线，结构如下：

\begin{center}
\small
\begin{tikzpicture}[
  node distance=5mm and 7mm, >=Stealth,
  blk/.style={draw, rounded corners, align=center, font=\small, inner sep=3pt, fill=main!6, draw=main!60!black},
  grp/.style={draw, rounded corners, align=center, font=\small, inner sep=3pt, fill=second!6, draw=second!60!black},
  ln/.style={->, thick, gray!60!black}]
\node[blk] (zmp) {ZMP/捕获点\\（\cref{ch:legged-biped} 地基）};
\node[blk, right=of zmp] (cmp) {CMP\\\,$=$ZMP$+\frac{1}{f_z}[\tau_{G,y};-\tau_{G,x}]$};
\node[blk, right=of cmp] (dash) {三级恢复仪表盘\\踝/髋臂/迈步};
\node[grp, below=8mm of zmp] (fly) {飞轮模型\\$\boldsymbol\xi_{eq}$ 被角动量拉近};
\node[grp, right=of fly] (cmm) {CMM $\mathbf h_G{=}\mathbf A_G\mathbf v$\\（\cref{ch:unified-multimodal-mpc}）};
\node[grp, right=of cmm] (task) {角动量任务\\$\mathbf A_G^{\mathrm{ang}}\dot{\mathbf v}{=}\dot{\mathbf k}_G^d{-}\dot{\mathbf A}_G^{\mathrm{ang}}\mathbf v$};
\node[blk, below=8mm of fly, fill=third!8, draw=third!60!black] (wbc) {TSID/HQP 任务栈\\CoM+角动量+摆脚+姿态};
\node[blk, right=of wbc, fill=third!8, draw=third!60!black] (ds) {双支撑过渡\\ZMP 转移+力 ramp};
\draw[ln] (zmp)--(cmp); \draw[ln] (cmp)--(dash);
\draw[ln] (cmp.south) -- ++(0,-4mm) -| (fly.north);
\draw[ln] (fly)--(cmm); \draw[ln] (cmm)--(task);
\draw[ln] (task.south) -- ++(0,-4mm) -| (wbc.north);
\draw[ln] (wbc)--(ds);
\draw[ln] (dash.south) to[out=-90,in=0] (task.east);
\end{tikzpicture}
\end{center}

\noindent\textbf{推荐路径}：\cref{sec:hmw-cmp} 用 CMP 把“角动量到底把平衡点移到哪”说精确，并由此给出踝/髋臂/迈步的几何仪表盘，是全章的概念起点，任何读者都该读；\cref{sec:hmw-flywheel} 是它的定量补充（角动量买来多少“等效捕获点”的余量）；\cref{sec:hmw-momentum-wbc} 是把这套思想\emph{落进全身 QP}的核心——角动量任务的推导与人形任务栈装配，写控制器前必须吃透；\cref{sec:hmw-ds} 是“两脚交接的那一瞬间怎么不抖”的工程细节，调试过渡期力矩跳变时回来看。

\subsection*{前置知识桥接}
本章把三块前置直接当作工具，凡前面已更深处讲过的，本章只在用到处 交叉引用 回引、不复述：
\begin{itemize}
  \item \cref{ch:legged-biped}（欠驱点足）给了本章的平衡语言地基：LIP 动力学 $\ddot x=\omega^2(x-p)$（\cref{eq:lb-lip}）、发散分量 DCM $\boldsymbol\xi=\mathbf x+\dot{\mathbf x}/\omega$（\cref{eq:lb-dcm-def}）、其一维动力学 $\dot{\boldsymbol\xi}=\omega(\boldsymbol\xi-\mathbf p)$（\cref{eq:lb-dcm-dyn}）与捕获点 $\mathbf p=\boldsymbol\xi$（\cref{eq:lb-capture}）。本章承接这些\emph{不再重证}，只把“迈步是唯一手段”放宽为“踝/髋臂/迈步三级”。其 ZMP/DCM 谱系亦见 \cref{sec:rm-dcm}（\cref{eq:rm-dcm-def}）；
  \item \cref{ch:unified-multimodal-mpc}（统一动力学与多模态 MPC）给了 CMM $\mathbf h_G=\mathbf A_G(\mathbf q)\mathbf v$（\cref{sec:umm-cmm}，\cref{eq:umm-cmm-def}，本书取 $\mathbf h_G=[\mathbf l_G;\mathbf k_G]$ 线在前）与“重力不产生绕质心力矩 $\Rightarrow$ 角动量守恒 $\dot{\mathbf k}_G=\mathbf 0$”（\cref{sec:umm-momentum}）。本章把 CMM 的\emph{角行} $\mathbf A_G^{\mathrm{ang}}$ 拿来组装角动量任务，把守恒律读成“臂动则腿/躯干反向动”的物理（臂作反作用轮），故不重述 CMM 定义；CMM 亦见 \cref{eq:rm-cmm}；
  \item \cref{ch:force-wbc}（力控与 WBC）给了 WBC 的统一动力学硬约束与 WBC-QP 标准型（\cref{eq:fw-wbcqp}）、“为何需要 WBC”的立论（\cref{sec:fw-why-wbc}）与笛卡尔阻抗（\cref{thm:fw-cart-impedance}）。本章的角动量任务、CoM 任务都\emph{插进}这个 QP；四足侧的任务优先级、零空间投影与加权/分层 QP 见 \cref{ch:quad_wbc}（\cref{sec:qw-priority}、\cref{sec:qw-nullspace}、\cref{eq:qw-qp-standard}、\cref{eq:qw-wqp}），本章只在“人形比四足多一条角动量任务”处对照。
\end{itemize}

\subsection*{如果跳过本章会怎样}
\begin{itemize}
  \item 你会把 \cref{ch:legged-biped} 的“平衡只能迈步”当成全部真相，从而\emph{看不见}脚掌带来的踝/髋臂两档余量——于是在小扰动下也急着迈步，步态频繁、僵硬；更糟的是，你不知道“摆臂救场”在数学上对应 CMP 出脚、对应一条可以塞进 QP 的角动量任务，于是人形受侧推时上体角动量在躯干里默默累积、直到 ZMP 出界翻倒，你却以为是落脚没规划好（这正是 \cref{ins:hmw-no-momentum-task} 要纠正的）。
  \item 你会把双支撑相当成“两脚都硬撑、随便分力”的过渡，从而在相位切换处算出阶跃的 ZMP 参考、阶跃的接触力、阶跃的关节力矩——实机上每步“咯噔”一下，重则激起结构振动，而你查遍单支撑公式都找不到根（根在 \cref{sec:hmw-ds} 的等效 $\omega_0$ 与力 ramp）。
\end{itemize}

\begin{table}[H]\centering\small
\caption{本章阅读方式建议}
\begin{tabular}{lll}
\toprule
阅读方式 & 时间 & 适合谁 \\
\midrule
精读（含 CMP 推导、角动量任务推导与练习） & 3--4 小时 & 首次系统学习人形动量 WBC \\
速读（跳过 ALIP/飞轮定量与 DS 多项式细节） & 1 小时 & 已会点足落脚平衡、想补人形角动量 \\
速查（只看小结的符号表 + 公式速查表） & 15 分钟 & 写 WBC 代码时回来查 CMP/角动量任务 \\
\bottomrule
\end{tabular}
\end{table}

\begin{note}
\textbf{本章记号与约定.}\;
沿用全书：向量粗体小写、矩阵粗体大写；旋转 $\mathbf R\in\SO(3)$、反对称算子 $(\cdot)^\wedge$（\cref{ch:lie}）；右扰动。
\textbf{质心量}：质心位置 $\mathbf r_G$（分量 $x,y,z$）、质心高度 $z_c$（恒值，LIP 假设）、总质量 $m$、重力加速度 $g$、LIP 自然频率 $\omega_0=\sqrt{g/z_c}$（这与 \cref{ch:legged-biped} 用 $\omega$、高度记 $h$ 是同一量，本章随源记号写 $\omega_0$、$z_c$）。
\textbf{动量}：质心空间动量 $\mathbf h_G=[\mathbf l_G;\mathbf k_G]$（线在前，\cref{eq:umm-cmm-def}），线动量 $\mathbf l_G=m\dot{\mathbf r}_G$、绕质心角动量 $\mathbf k_G$；其变化率 $\dot{\mathbf k_G}$ 即绕质心力矩 $\boldsymbol\tau_G$。CMM 记 $\mathbf A_G$，角行块记 $\mathbf A_G^{\mathrm{ang}}$（后三行）。
\textbf{平衡点}：ZMP（零力矩点）$\mathbf r_{\mathrm{ZMP}}$、CoP（压心）、CMP（质心力矩枢轴点）$\mathbf r_{\mathrm{CMP}}$、DCM/捕获点 $\boldsymbol\xi$（同 \cref{eq:lb-dcm-def}）。地面法向力 $f_z$。
\textbf{动力学}：浮动基统一动力学 $\mathbf M(\mathbf q)\dot{\mathbf v}+\mathbf h=\mathbf S^\top\boldsymbol\tau+\sum_c\mathbf J_c^\top\boldsymbol\lambda$（\cref{eq:fw-wbcqp}），其中 $\mathbf S$ 选驱动行、$\mathbf J_c$ 接触 Jacobian、$\boldsymbol\lambda$ 接触力；末端力到关节力矩 $\boldsymbol\tau=\mathbf J^\top\mathbf f$。
\textbf{步态量}：步宽 $W$、步时 $T$、双支撑相 DS、过渡相位标量 $s:0\!\to\!1$。
\end{note}

% ════════════════════════════════════════════════════════════════
\section{CMP 与三级恢复策略仪表盘}
\label{sec:hmw-cmp}

\paragraph{动机：当机器人动用角动量，地面反力的等效作用线落到哪？}
\cref{ch:legged-biped} 反复强调一条铁律：ZMP 必须落在支撑域内，否则脚底要“拉地面”才能实现等效力矩，物理上不可能。这条铁律建立在一个被悄悄假设的前提上——\textbf{角动量不变}（$\dot{\mathbf k}_G\approx\mathbf 0$，正是 LIPM 的核心假设，见 \cref{ins:hmw-lipm-assumption}）。可一旦机器人主动甩臂、扭腰，绕质心就有了非零力矩 $\dot{\mathbf k}_G\ne\mathbf 0$，这条前提破了。于是出现一个新问题：算上角动量这部分效应后，\emph{地面反力的合力作用线}与地面的交点落在哪？这个点就是 CMP（Centroidal Moment Pivot，质心力矩枢轴点）。它正是“摆臂救平衡”这件事在几何上的精确刻画。

\begin{insight}[LIPM 把角动量项扔掉了，CMP 是把它捡回来]\label{ins:hmw-lipm-assumption}
回忆 \cref{eq:lb-lip} 的 LIP 动力学 $\ddot x=\omega_0^2(x-p)$，它干净利落地把 ZMP $p$ 与质心加速度直接挂钩。但这份干净是有代价的——它假设 $\dot{\mathbf k}_G=\mathbf 0$，即\emph{忽略一切绕质心的角动量变化}。完整的 ZMP 公式（从质心角动量定理出发）其实带一个角动量项：例如矢状方向 $p_x=x-\dfrac{z_c\ddot x}{g}-\dfrac{\dot k_{G,y}}{m\,g}$（取 $\ddot z=0$；此负号与 Popović eq.~(6) 及 \cref{eq:hmw-cmp-def} 一致），那个 $\dot k_{G,y}/(mg)$ 的量纲恰是\emph{长度}。\textbf{LIPM 是把这一长度项强行置零}。CMP 不做这个置零——它正是回答“把这部分角动量效应折算成长度后，等效压力点移到了哪”。所以 CMP 与 ZMP 的差，量化的就是“你现在动用了多少角动量”。
\end{insight}

\subsection{CMP 的定义与它和 ZMP 的关系}
\label{subsec:hmw-cmp-def}

把绕质心的力矩（即角动量变化率）记作 $\boldsymbol\tau_G=\dot{\mathbf k}_G$。CMP 定义为\emph{地面反力作用线与地面的交点}；当系统动用角动量产生 $\boldsymbol\tau_G$ 时，这条作用线相对“无角动量时的 ZMP”发生水平偏移，二者关系为

\begin{equation}
  \boxed{\;\mathbf r_{\mathrm{CMP}}=\mathbf r_{\mathrm{ZMP}}+\frac{1}{f_z}\begin{bmatrix}\tau_{G,y}\\[2pt]-\tau_{G,x}\end{bmatrix}\;}
  \label{eq:hmw-cmp-def}
\end{equation}

其中 $f_z$ 是地面法向力，$\tau_{G,x},\tau_{G,y}$ 是绕质心力矩 $\boldsymbol\tau_G=\dot{\mathbf k}_G$ 的水平两分量。此式即 Popović–Goswami–Herr 的 CMP 原始定义（\cite{popovic2005cmp}，IJRR 24(12):1013–1032，2005，eq.~(20)：$x_{\mathrm{CMP}}=x_{\mathrm{ZMP}}+\tau_y/F_z$、$y_{\mathrm{CMP}}=y_{\mathrm{ZMP}}-\tau_x/F_z$），\emph{分母确为地面法向力 $F_z=f_z$ 而非 $mg$}（仅当 $\ddot z=0$ 时 $f_z=mg$ 二者数值相等）。由质心角动量定理 $\dot{\mathbf k}_G=(\mathbf r_{\mathrm{ZMP}}-\mathbf r_G)\times\mathbf F$ 直接展开可独立证得上式正负号（即 $r_{\mathrm{CMP},x}-r_{\mathrm{ZMP},x}=+\tau_{G,y}/f_z$、$r_{\mathrm{CMP},y}-r_{\mathrm{ZMP},y}=-\tau_{G,x}/f_z$）。直觉上，分子的 $\boldsymbol\tau_G$ 是“角动量贡献的力矩”，除以法向力 $f_z$ 得到一段\emph{长度}——它就是 CMP 相对 ZMP 的水平位移。当 $\dot{\mathbf k}_G=\mathbf 0$（不动用角动量）时，\cref{eq:hmw-cmp-def} 给出 $\mathbf r_{\mathrm{CMP}}=\mathbf r_{\mathrm{ZMP}}$，两点重合——这正是 LIPM 世界里的情形。

\begin{insight}[一句话分清三个点：ZMP 必在脚内，CMP 可出脚，捕获点指迈步]\label{ins:hmw-three-points}
\cref{ch:legged-biped} 反复用 ZMP 与捕获点 $\boldsymbol\xi$，本章补上中间的 CMP。三者回答的是\emph{三个不同问题}：
\begin{itemize}
  \item \textbf{ZMP} 问“地面反力的等效作用点在哪、它\emph{必须}在脚掌内”——因为脚底只能压、不能拉。它是\emph{约束}。
  \item \textbf{CMP} 问“算上角动量后，等效作用线落在哪、它\emph{可以}在脚掌外”——脚外的那部分，正是用角动量“虚拟地”买来的力矩。它是\emph{状态读数}。
  \item \textbf{捕获点} $\boldsymbol\xi$ 问“我该往哪踩才能停下”（\cref{eq:lb-capture}）。它是\emph{动作目标}。
\end{itemize}
关键在于：CMP 与 ZMP 的差告诉你“正在动用多少角动量”，捕获点是否出脚告诉你“是否非迈步不可”。把这三点同画在一张俯视图上（\cref{fig:hmw-three-points}），就是一台完整的平衡仪表盘。
\end{insight}

\subsection{踝、髋/臂、迈步：三级恢复策略的几何区分}
\label{subsec:hmw-three-strategies}

CMP 的教学价值，正在于它把人（和人形）应对扰动的三档反应\emph{几何化}了。它们的区别不是“用哪个关节”这么模糊，而是“CMP 与 ZMP 的相对位置、捕获点是否出脚”这么精确——见 \cref{tab:hmw-strategies}。

\begin{table}[H]\centering\small
\caption{三级恢复策略的几何区分（按扰动从小到大）}
\begin{tabular}{p{0.20\textwidth} p{0.26\textwidth} p{0.22\textwidth} p{0.18\textwidth}}
\toprule
策略 & 几何特征 & 物理动作 & 何时用 \\
\midrule
踝策略 (ankle) & CMP $=$ ZMP，都在脚内 & 只靠踝关节力矩在脚掌内移动 CoP/ZMP & 小扰动 \\
髋/臂策略 (hip) & CMP 跑到脚\emph{外}，ZMP 仍在脚\emph{内} & 甩臂/扭腰产生 $\dot{\mathbf k}_G\ne\mathbf 0$ & 中大扰动 \\
迈步 (stepping) & 连 CMP 策略也不够、捕获点出脚 & 改变支撑（迈一步，\cref{eq:lb-capture}） & 大扰动 \\
\bottomrule
\end{tabular}
\label{tab:hmw-strategies}
\end{table}

这张表的逻辑链是单调递进的：扰动小，踝部足以把 ZMP 拉到平衡所需处，此时不必动角动量，CMP 与 ZMP 重合；扰动变大，ZMP 已经顶到脚掌边缘还不够，机器人开始甩臂——角动量项把\emph{等效}作用点（CMP）推到脚掌外，而真实的 ZMP 仍被钉在脚内边缘（它\emph{不能}出脚），这一“CMP 出脚、ZMP 不出脚”的劈叉，正是髋/臂策略的几何指纹；扰动再大，连角动量也撑不住（下节说明它会饱和），捕获点已经跑出脚外，只剩迈步一条路（回到 \cref{ch:legged-biped} 的捕获点落脚）。

\begin{figure}[ht]
\centering
\begin{tikzpicture}[scale=1.0,>=Stealth]
  % three foot rectangles (top view), one per strategy
  \foreach \dx/\lab in {0/{踝策略}, 4.2/{髋臂策略}, 8.4/{迈步}}{
    \begin{scope}[shift={(\dx,0)}]
      % foot rectangle
      \draw[gray!70!black,thick,fill=gray!6] (-0.55,-1.0) rectangle (0.55,1.0);
      \node[gray!55!black] at (0,1.35) {\small \lab};
    \end{scope}
  }
  % --- ankle: CMP=ZMP inside ---
  \begin{scope}[shift={(0,0)}]
    \fill[main!70!black] (0.15,0.2) circle (2.2pt);
    \node[main!70!black,font=\scriptsize] at (0.15,0.55) {ZMP=CMP};
    \node[align=center,font=\scriptsize,gray!55!black] at (0,-1.45) {同点、都在脚内\\$\dot{\mathbf k}_G=\mathbf 0$};
  \end{scope}
  % --- hip/arm: CMP out, ZMP at edge inside ---
  \begin{scope}[shift={(4.2,0)}]
    \fill[main!70!black] (0.5,0.1) circle (2.2pt) node[above,font=\scriptsize,main!70!black]{ZMP};
    \fill[second!60!black] (1.45,0.1) circle (2.2pt) node[above,font=\scriptsize,second!60!black]{CMP};
    \draw[second!60!black,dashed,->] (0.5,0.1)--(1.4,0.1);
    \node[second!60!black,font=\scriptsize] at (0.95,-0.25) {$\frac{1}{f_z}\boldsymbol\tau_G$};
    \node[align=center,font=\scriptsize,gray!55!black] at (0.3,-1.45) {CMP 出脚\\ZMP 抵边、不出脚};
  \end{scope}
  % --- stepping: capture point out ---
  \begin{scope}[shift={(8.4,0)}]
    \fill[main!70!black] (0.5,0.0) circle (2.2pt) node[above,font=\scriptsize,main!70!black]{ZMP};
    \fill[teal!55!black] (1.7,0.0) circle (2.4pt) node[above,font=\scriptsize,teal!55!black]{$\boldsymbol\xi$};
    \draw[teal!55!black,->,thick] (0.6,0.0)--(1.6,0.0);
    % new footstep
    \draw[third,thick,fill=third!8] (1.25,-0.9) rectangle (2.35,0.6);
    \node[third,font=\scriptsize] at (1.8,-1.15) {新落脚};
    \node[align=center,font=\scriptsize,gray!55!black] at (1.0,-1.55) {捕获点出脚\\必须迈步};
  \end{scope}
\end{tikzpicture}
\caption{三级恢复策略的平衡仪表盘（脚掌俯视）。左：踝策略，CMP 与 ZMP 同点、都在脚内，不动用角动量。中：髋/臂策略，甩臂产生 $\dot{\mathbf k}_G\ne\mathbf 0$，CMP 被 \cref{eq:hmw-cmp-def} 的 $\frac{1}{f_z}\boldsymbol\tau_G$ 推出脚外，而真实 ZMP 仍被钉在脚掌边缘内（它不能出脚）——“CMP 出脚、ZMP 抵边”就是动用角动量的几何指纹。右：连角动量也不够，捕获点 $\boldsymbol\xi$（\cref{eq:lb-capture}）跑出脚外，只能迈步、把支撑移到新落脚处。读 CMP 与 ZMP 之差、读 $\boldsymbol\xi$ 是否出脚，就读出了“机器人正在用哪一档”。}
\label{fig:hmw-three-points}
\end{figure}

\begin{insight}[这正是 \cref{ch:legged-biped} “只能迈步”的放宽：脚掌买来了两档]\label{ins:hmw-relax-stepping}
\cref{ins:lb-must-step} 在点足上证明了“纯反馈救不回、唯一手段是迈步”，根因是\emph{点接触}：踝部无力矩、CoP 不能在“脚掌”内移动（点足根本没有脚掌），支撑点一步内钉死。本章的人形有\emph{有面积的脚掌}，于是那条结论被放宽——CoP 可以在脚掌内移动（踝策略这一档回来了），上体可以甩臂（髋/臂策略这一档也回来了），只有当这两档都顶满、捕获点仍出脚时，才退回 \cref{ch:legged-biped} 的迈步。\textbf{“点足只有迈步”是“人形三级策略”在脚掌面积$\to 0$时的退化极限。}理解这一点，你就知道 ZMP/CMP/捕获点不是三个孤立概念，而是同一根“支撑能力”坐标轴上从弱到强的三个刻度。
\end{insight}

\begin{pitfall}[误区：以为“CMP 能出脚”就等于角动量无限可用]\label{pit:hmw-cmp-saturation}
CMP 可以跑到脚掌外，很容易让人以为“那就多甩臂、把 CMP 推得远远的，平衡能力就无限了”。这是错的，而且错得危险。CMP 在脚外\emph{是有代价的}——它对应着角动量在\emph{持续累积}（$\dot{\mathbf k}_G\ne\mathbf 0$ 一直为同号），而角动量\emph{不能无限增大}：躯干能扭的角度有限、手臂的行程有限。一旦臂摆到极限、躯干扭到极限，就再没有 $\dot{\mathbf k}_G$ 可产生，CMP 被迫缩回脚内，平衡能力骤降。\textbf{用角动量买来的是\emph{时间}，不是无限的平衡能力。}CMP 长时间停在脚外，意味着角动量很快会饱和，之后仍必须迈步。下一节的飞轮模型（\cref{sec:hmw-flywheel}）把这个“买时间”说成定量的——它精确地撑大了多少“不迈步的可恢复区”。
\end{pitfall}

\begin{pitfall}[误区：把 ZMP 与 CoP 当成一回事]\label{pit:hmw-zmp-cop}
ZMP 与 CoP（压心）在\emph{平地、忽略足底摩擦力矩}时数值一致，于是常被混为一谈。但它们的\emph{身份}不同：CoP 是一个纯力学测量量（由脚底力/力矩传感器读出的压力分布中心），ZMP 是一个动力学概念（由质心运动方程定义、控制器\emph{计算}出来）。在倾斜地面或存在足底扭矩时两者可不同。调试时应\emph{分别}记录——CoP 来自传感器、ZMP 来自模型，它们的差是“接触模型准不准”的体检指标。在仿真里两者通常一致（接触模型精确），到实机就未必，这一差异本身就是诊断接触建模误差的工具。
\end{pitfall}

\begin{practice}
\begin{enumerate}
  \item 由 \cref{eq:hmw-cmp-def} 出发，令 $\dot{\mathbf k}_G=\mathbf 0$ 验证 CMP 与 ZMP 重合；再设人形以 $\dot k_{G,y}=5\,\mathrm{N\!\cdot\!m}$ 的速率绕 $y$ 轴摆臂、$f_z=mg$、$m=40\,\mathrm{kg}$，算出 CMP 相对 ZMP 的矢状偏移。这段偏移相对一个 $12\,\mathrm{cm}$ 宽的脚掌有多大？
  \item 把 \cref{tab:hmw-strategies} 三行与 \cref{fig:hmw-three-points} 三图一一对应，用“CMP 与 ZMP 是否重合 / 捕获点是否出脚”两个布尔量，写出一个判别“当前是第几级策略”的决策表（提示：两个布尔量给出四种组合，其中一种不会出现，是哪一种、为什么）。
  \item （概念）\cref{ins:hmw-relax-stepping} 说点足是人形三级策略在“脚掌面积$\to 0$”的退化。请反过来想：若脚掌\emph{无限大}（机器人站在一整块和它一样大的平板上），踝策略的可调范围有多大？这时还需要髋/臂与迈步吗？这说明三级策略的“档位边界”由什么决定。
\end{enumerate}
\end{practice}

% ════════════════════════════════════════════════════════════════
\section{飞轮模型与等效捕获点}
\label{sec:hmw-flywheel}

\paragraph{动机：角动量到底把“不迈步的可恢复区”撑大了多少？}
\cref{sec:hmw-cmp} 定性地说“甩臂能买时间”。但工程上要的是定量答案：\emph{受到这么大的推、不迈步、靠甩臂，到底能不能稳住？} 回答它需要给 LIP 加一个能显式产生角动量的部件——\textbf{飞轮}（flywheel）。这正是 Pratt 等人提出捕获点框架时所用的模型（\cref{ch:legged-biped} 的 \cref{eq:lb-capture} 引用了同一脉络）。

\subsection{飞轮把等效捕获点拉近}
\label{subsec:hmw-flywheel-eq}

把 \cref{ch:legged-biped} 的“质点 + 无质量腿”倒立摆，升级为“质点 + 无质量腿 + 一个可绕质心转动的飞轮”。飞轮转动产生绕质心的角动量 $\Delta\mathbf k_G$，它在不移动支撑点的前提下，临时改变了质心的有效受力，从而把\emph{等效}捕获点拉近：

\begin{equation}
  \boxed{\;\boldsymbol\xi_{eq}=\boldsymbol\xi-\frac{\Delta\mathbf k_G}{m\,\omega_0\,z_c}\;}
  \label{eq:hmw-flywheel-xi}
\end{equation}

其中 $\boldsymbol\xi$ 是不动用角动量时的捕获点（\cref{eq:lb-capture}），$\Delta\mathbf k_G$ 是飞轮（躯干转动、摆臂）提供的\emph{瞬时}角动量（即 $\Delta\mathbf k_G=J\,\Delta\dot\theta$，飞轮转动惯量 $J$ 乘其角速度增量，非积分量），$m$ 为总质量、$\omega_0=\sqrt{g/z_c}$、$z_c$ 为质心高度。此式可由 Pratt 等的“LIP + 飞轮”模型直接得到\cite{pratt2006capture}：捕获点 $\xi=\dot x/\omega_0$（\cref{eq:lb-capture}），而飞轮一次冲量使质心有效水平速度产生阶跃 $\Delta\dot x=-J\,\Delta\dot\theta/(m z_0)=-\Delta\mathbf k_G/(m z_c)$（Pratt 原文 eq.~(11) 邻文），除以 $\omega_0$ 即得 \cref{eq:hmw-flywheel-xi}。\textbf{故分母确为 $m\,\omega_0\,z_c$（一次方 $\omega_0$、一次方 $z_c$），既非 $m\omega_0^2 z_c$ 亦非仅 $m\omega_0$}（联网核对 Pratt 2006 与 Koolen 2012 可捕获性框架\cite{koolen2012capturability} 均吻合；量纲亦自洽：$\Delta k_G$ 单位 $\mathrm{kg\,m^2/s}$，除以 $m\,\omega_0\,z_c$ 的 $\mathrm{kg\cdot s^{-1}\cdot m}$ 得长度 $\mathrm m$）。式子的方向很直观：合适符号的 $\Delta\mathbf k_G$ 让 $\boldsymbol\xi_{eq}$ 比 $\boldsymbol\xi$ 更靠近支撑脚——也就是把“本来已经跑出脚外、非迈步不可的捕获点”，临时\emph{拽回}脚内。一旦 $\boldsymbol\xi_{eq}$ 重新落进脚掌，机器人就能不迈步地稳住（至少暂时）。

\begin{insight}[等效捕获点 = 把“0 步可捕获区”暂时撑大]\label{ins:hmw-flywheel-region}
\cref{eq:hmw-flywheel-xi} 给“甩臂买时间”一个精确含义。\cref{ch:legged-biped} 的判据是二元的：捕获点在脚内则不迈步、在外则迈步。飞轮把这条边界\emph{暂时外推}：原本捕获点 $\boldsymbol\xi$ 已在脚外（按裸判据须迈步），但只要飞轮能提供足够的 $\Delta\mathbf k_G$ 把 $\boldsymbol\xi_{eq}$ 拉回脚内，机器人就\emph{暂时}不必迈步——等价于把“不迈步就能恢复”的区域撑大了一圈。这就是人在窄独木桥上被推时会\emph{猛地}转动手臂的物理原因：用一次性的、幅度很大的角动量，把瞬时捕获点硬拉回脚下，避免被迫迈步（在独木桥上迈步往往无处可迈）。\textbf{角动量买的是“暂时”二字——因为 $\Delta\mathbf k_G$ 受 \cref{pit:hmw-cmp-saturation} 的饱和所限，撑大的区域不能维持，臂摆到头就得还回去。}
\end{insight}

\subsection{角动量守恒视角：臂动则身反向动}
\label{subsec:hmw-conservation}

飞轮“凭空”产生角动量了吗？没有。\cref{sec:umm-momentum} 给的角动量守恒 $\dot{\mathbf k}_G=\mathbf 0$（重力不产生绕质心力矩，\cref{eq:umm-kG-cons}）说的是\emph{在腾空相、无外接触力矩时}总角动量守恒；而站立平衡时，脚与地面的接触\emph{能}提供外力矩，于是 $\dot{\mathbf k}_G$ 可以非零。这两件事不矛盾：

\begin{insight}[守恒与“摆臂救场”不矛盾：接触力矩是角动量的进出口]\label{ins:hmw-conservation-noncontradict}
前置自测第 3 题问的正是这个表面矛盾。关键在于\emph{是否有外接触力矩}：
\begin{itemize}
  \item \textbf{腾空 / 无外力矩}时，唯一外力是重力，重力线过质心不产生力矩，故 $\dot{\mathbf k}_G=\mathbf 0$、总角动量守恒（\cref{eq:umm-kG-cons}）。此时“摆臂”只能在\emph{内部重新分配}角动量——臂朝一个方向转，身体/腿必朝反方向转以保持总量不变（这正是猫空中翻身、跳水转体的原理，臂/腿互为反作用轮）。\cref{sec:umm-cmm} 用 CMM 的角行块 $\mathbf A_G^{\mathrm{ang}}$ 精确描述了这种“谁动多少、总和为常数”的分配。
  \item \textbf{站立、有接触}时，脚底的地面反力相对质心有力臂，能产生绕质心的外力矩 $\boldsymbol\tau_G=\dot{\mathbf k}_G\ne\mathbf 0$。此时甩臂\emph{真能}改变总角动量——多出来/少掉的那部分，正是经由脚底接触力矩与地面交换的。CMP 出脚（\cref{eq:hmw-cmp-def}）量化的就是这笔交换。
\end{itemize}
所以“守恒”和“摆臂救场”分属两个相位，靠“有没有接触力矩”这道闸门分开。理解这道闸门，你才知道为什么\emph{腾空相}不能靠摆臂改变落点动量（只能内部调姿态），而\emph{支撑相}可以。
\end{insight}

\begin{note}[臂作反作用轮：与 \cref{ch:unified-multimodal-mpc} 去重]\label{note:hmw-arm-reaction-wheel}
“臂朝一方向转、身体反向转以守恒总角动量”这件事，\cref{sec:umm-momentum} 已在多模态相位的角动量守恒里讲过（臂作反作用轮、CCRBA/CMM 视角）。本章\emph{不重复}那套推导，只把它接到双足平衡场景：飞轮模型 \cref{eq:hmw-flywheel-xi} 里的 $\Delta\mathbf k_G$，在站立相是经接触力矩从地面“借”来的、可用于救平衡；在腾空相则只能内部分配、不能改变落点。需要 CMM 角行块如何把全身关节速度映射到 $\mathbf k_G$ 的细节，回到 \cref{sec:umm-cmm}、\cref{eq:umm-cmm-def}。
\end{note}

\begin{practice}
\begin{enumerate}
  \item 在 \cref{eq:hmw-flywheel-xi} 中取 $z_c=0.8\,\mathrm{m}$、$\omega_0=\sqrt{9.81/0.8}\approx 3.5\,\mathrm{rad/s}$、$m=40\,\mathrm{kg}$。设裸捕获点 $\boldsymbol\xi$ 在脚掌外缘外 $3\,\mathrm{cm}$，问飞轮需提供多大的 $\Delta\mathbf k_G$ 才能把 $\boldsymbol\xi_{eq}$ 拉回外缘（即拉近 $3\,\mathrm{cm}$）？（\cref{eq:hmw-flywheel-xi} 的分母 $m\omega_0 z_c$ 已对 Pratt 2006 原文核实；本题给出的是单飞轮极简模型下的量级估计，真实人形需经全身 QP 协调多关节产生该 $\Delta\mathbf k_G$。）
  \item 用 \cref{ins:hmw-conservation-noncontradict} 解释：跳水运动员在空中抱团转体很快、展开后转体变慢——这是 $\dot{\mathbf k}_G=\mathbf 0$ 下的什么效应？它和站立时“甩臂改变总角动量”有何本质不同？
  \item （承接 \cref{pit:hmw-cmp-saturation}）把飞轮的可用角动量上限记为 $\Delta k_{\max}$（受臂行程/躯干转角限）。由 \cref{eq:hmw-flywheel-xi} 写出“飞轮能把捕获点最多拉近多少”的表达式。这个上限说明了什么时候\emph{再怎么甩臂也必须迈步}。
\end{enumerate}
\end{practice}

% ════════════════════════════════════════════════════════════════
\section{人形动量 WBC：从 CMM 到角动量任务与任务栈}
\label{sec:hmw-momentum-wbc}

\paragraph{动机：把“管住角动量”落进全身 QP。}
\cref{sec:hmw-cmp,sec:hmw-flywheel} 在低维模型上说清了“角动量能救平衡”。但真实人形有 $30+$ 个关节，最终要解的是一个全身 QP（\cref{eq:fw-wbcqp}）——“甩臂”不是一个抽象指令，而是 QP 自己选出的、满足动力学与接触约束的关节加速度。要让 QP “主动用角动量救平衡”，就得在它的任务栈里\emph{显式}加一条“角动量任务”。本节把这条任务从 CMM 推出来，并装配人形完整任务栈。

\subsection{角动量任务：从 CMM 求导而来}
\label{subsec:hmw-momentum-task}

地基是 \cref{ch:unified-multimodal-mpc} 的 CMM（\cref{eq:umm-cmm-def}）：质心空间动量 $\mathbf h_G=\mathbf A_G(\mathbf q)\mathbf v$，本书取 $\mathbf h_G=[\mathbf l_G;\mathbf k_G]$（线在前），故 CMM 的\emph{后三行} $\mathbf A_G^{\mathrm{ang}}$ 把广义速度映射到角动量 $\mathbf k_G=\mathbf A_G^{\mathrm{ang}}\mathbf v$。WBC-QP 的决策变量含广义加速度 $\dot{\mathbf v}$（\cref{eq:fw-wbcqp}），故任务必须写成“关于 $\dot{\mathbf v}$ 的线性式”。对 $\mathbf h_G=\mathbf A_G\mathbf v$ 两边求时间导数：

\begin{equation}
  \dot{\mathbf h}_G=\mathbf A_G(\mathbf q)\,\dot{\mathbf v}+\dot{\mathbf A}_G(\mathbf q,\mathbf v)\,\mathbf v,
  \label{eq:hmw-hGdot}
\end{equation}

其中 $\dot{\mathbf A}_G$ 是 CMM 随构型变化的时间导数（数值上可由 Pinocchio \texttt{dccrba} 算得，与 \texttt{ccrba} 给出的 $\mathbf A_G$ 配套）。取角分量（后三行），并把它整理成“任务加速度 = 已知量”的标准形式：

\begin{equation}
  \boxed{\;\mathbf A_G^{\mathrm{ang}}\,\dot{\mathbf v}=\dot{\mathbf k}_G^{d}-\dot{\mathbf A}_G^{\mathrm{ang}}\,\mathbf v\;}
  \label{eq:hmw-angular-task}
\end{equation}

这就是\textbf{角动量任务}：Jacobian 是 CMM 的角行块 $\mathbf A_G^{\mathrm{ang}}$，期望任务加速度右端 $\mathbf b=\dot{\mathbf k}_G^{d}-\dot{\mathbf A}_G^{\mathrm{ang}}\mathbf v$。注意那个偏置项 $-\dot{\mathbf A}_G^{\mathrm{ang}}\mathbf v$ \emph{不能漏}——它是“速度依赖的偏置角动量变化率”，物理上对应当前运动本身带来的角动量漂移，与接触任务里不能漏 $\dot{\mathbf J}_c\mathbf v$（\cref{eq:fw-wbcqp}）是同一类项。

\paragraph{期望角动量变化率 $\dot{\mathbf k}_G^d$ 怎么给。}
最常见的三种取法，对应三种平衡意图：
\begin{itemize}
  \item \textbf{零角动量}（站立、缓步）：令角动量收敛到零，$\dot{\mathbf k}_G^d=-\mathbf K_k\,\mathbf k_G=-\mathbf K_k\,\mathbf A_G^{\mathrm{ang}}\mathbf v$。代回 \cref{eq:hmw-angular-task} 得 $\mathbf A_G^{\mathrm{ang}}\dot{\mathbf v}=-\mathbf K_k\mathbf A_G^{\mathrm{ang}}\mathbf v-\dot{\mathbf A}_G^{\mathrm{ang}}\mathbf v$。这把任意累积的角动量主动“泄掉”，是站立抗扰的默认。
  \item \textbf{恒定角动量}（惯性滑行）：$\dot{\mathbf k}_G^d=\mathbf 0$，保持当前角动量不变。
  \item \textbf{前馈角动量}（高动态、主动甩臂）：$\dot{\mathbf k}_G^d=\dot{\mathbf k}_G^{\mathrm{plan}}$，跟踪规划层给出的角动量曲线——这正是“主动用髋/臂策略”在 QP 里的表达，让飞轮模型 \cref{eq:hmw-flywheel-xi} 那笔“买时间”的角动量，真正由 QP 协调全身关节产生。
\end{itemize}
增益 $\mathbf K_k$ 取大则角动量被快速压平、但易激起躯干/臂高频抖动（要大关节加速度），取小则恢复慢、扰动期 DCM 可能已发散。

\begin{insight}[不加角动量任务，人形受侧推必“身越来越斜”直到 ZMP 出界]\label{ins:hmw-no-momentum-task}
四足 WBC（\cref{ch:quad_wbc}）的任务栈通常没有角动量任务，照样跑得稳，于是初学者很自然地在人形上也省掉它。在平地慢走时这几乎无害——角动量变化本就很小。但一旦人形受到侧向推力，差异是\emph{致命}的：
\begin{itemize}
  \item \textbf{有角动量任务}：QP 主动选择躯干、臂、腿的协同加速度去满足 $\dot{\mathbf k}_G^d$，躯干快速回正、臂做补偿动作把累积的角动量吸收掉（CMP 短暂出脚、随即收回，对应 \cref{fig:hmw-three-points} 中图）。
  \item \textbf{无角动量任务}：没有任何任务去约束 $\mathbf k_G$，角动量在上体里\emph{持续累积}、躯干\emph{持续}倾斜，CoM 任务只管线动量、管不到这一项，机器人一路斜下去直到 ZMP 顶出脚掌边界翻倒。
\end{itemize}
这正是 IHMC 在 DARPA 机器人挑战赛里把\emph{动量目标}放在最高优先级的原因——它是抗扰恢复的关键，也是“人形 WBC 必有、四足 WBC 可无”这条分界的根（练习见本节末）。\textbf{四足质量分布紧凑、躯干转动惯量大、单步冲击小，把躯干姿态做好就够；人形上体高、臂长，扰动下角动量在躯干快速累积，非显式管不可。}
\end{insight}

\subsection{人形 TSID/HQP 任务栈装配}
\label{subsec:hmw-task-stack}

把上述角动量任务，连同 CoM、摆脚、姿态等任务，按优先级装进 WBC-QP（\cref{eq:fw-wbcqp}），就是人形动量 WBC（IHMC/TALOS 风格）。决策变量取广义加速度与接触力 $[\dot{\mathbf v};\boldsymbol\lambda]$，关节力矩由动力学方程反推

\begin{equation}
  \boldsymbol\tau=\mathbf S\big(\mathbf M\dot{\mathbf v}+\mathbf h-\textstyle\sum_c\mathbf J_c^\top\boldsymbol\lambda\big),
  \label{eq:hmw-torque-recover}
\end{equation}

即 \cref{ch:force-wbc} 的 AccForce 思路（消去 $\boldsymbol\tau$ 缩小 QP）。硬约束与四足同构：浮动基不可驱动行 $\mathbf S_f(\mathbf M\dot{\mathbf v}+\mathbf h-\sum_c\mathbf J_c^\top\boldsymbol\lambda)=\mathbf 0$、支撑脚不滑 $\mathbf J_c\dot{\mathbf v}+\dot{\mathbf J}_c\mathbf v=\mathbf 0$、摩擦锥（线性化金字塔）与脚掌 CoP 边界（\cref{ch:force-wbc} 已详述，本章不重复）。任务栈见 \cref{tab:hmw-task-stack}。

\begin{table}[H]\centering\small
\caption{人形动量 WBC 任务栈（优先级从高到低；对照四足 \cref{ch:quad_wbc}）}
\begin{tabular}{p{0.16\textwidth} p{0.18\textwidth} p{0.24\textwidth} p{0.16\textwidth}}
\toprule
任务 & Jacobian & 参考 & 优先级 \\
\midrule
支撑接触（不滑） & $\mathbf J_c$ & 零接触加速度 & 硬约束 \\
浮动基动力学 & $\mathbf S_f(\cdots)$ & $=\mathbf 0$ & 硬约束 \\
CoM / 线动量 & $\mathbf J_{\mathrm{com}}$（即 $\mathbf A_G^{\mathrm{lin}}/m$） & LIPM/DCM 参考（\cref{eq:lb-dcm-def}） & 高 \\
\textbf{角动量（人形新增）} & $\boldsymbol{A}_G^{\mathrm{ang}}$ & $\dot{\mathbf k}_G^d$（\cref{eq:hmw-angular-task}） & 高 \\
摆动脚 & $\mathbf J_{\mathrm{swing}}$ & 落脚轨迹（\cref{sec:hmw-ds}） & 中高 \\
关节姿态 & $\mathbf I$ & 默认姿态 & 低 \\
\bottomrule
\end{tabular}
\label{tab:hmw-task-stack}
\end{table}

\begin{insight}[人形 vs 四足任务栈：差的就是那条角动量任务]\label{ins:hmw-vs-quad}
把 \cref{tab:hmw-task-stack} 与四足 WBC/WBIC（\cref{ch:quad_wbc}，\cref{eq:qw-qp-standard}/\cref{eq:qw-wqp}）逐行对照，会发现支撑接触、浮动基动力学、CoM、摆脚、姿态\emph{几乎一一对应}——真正多出来的，是高亮那一行\textbf{角动量任务}。这不是偶然，而是 \cref{ins:hmw-no-momentum-task} 的结构性后果。优先级排布上，IHMC 一系把\emph{动量}（线 + 角合成的 $\dot{\mathbf h}_G$）当作最高级优化目标，因为动量直接通过牛顿—欧拉关系 $\dot{\mathbf h}_G=\sum_c\mathbf f_c+m\mathbf g$ 关联接触力——控制动量等价于控制接触力的合效果，且 $6$ 维动量目标的 QP 规模很小。\cref{ch:quad_wbc} 的零空间投影（\cref{sec:qw-nullspace}）与字典序优先级（\cref{sec:qw-priority}）机制在人形上原样适用，只是被投影/排序的任务多了角动量这一条。
\end{insight}

\begin{insight}[HQP 之于角动量：把它放进高层零空间，低层不准动它]\label{ins:hmw-hqp-momentum}
为什么角动量任务要放\emph{高}优先级、而非和操作任务平权地加权？这正是 \cref{ch:quad_wbc} “加权 QP vs HQP”之争（\cref{eq:qw-wqp}）在人形上的关键落点。若把角动量任务与“双臂末端跟踪”同权加权，当臂去抓一个远处重物时，臂任务的大 Jacobian 会让加权残差偏向迁就臂、\emph{牺牲}角动量精度——优先级泄漏，躯干悄悄倾斜直到 ZMP 出界。HQP 用零空间投影（\cref{sec:qw-nullspace}）把高优先级任务的最优值\emph{冻结}进下层约束：角动量与 CoM 进高层，操作任务只能在“不改变平衡”的零空间里腾挪。\textbf{对人形，平衡（CoM + 角动量）必须在操作之上，且要用 HQP 的硬投影而非加权的软偏好来保证——只要边界会被触碰，软约束就不够。}
\end{insight}

\begin{note}[ZMP 裕度 / DCM 误差 / 角动量峰值：天然的奖励信号（连 P8 学习控制）]\label{note:hmw-reward-signals}
本章这一整套有物理含义的标量——ZMP 相对脚掌边界的\emph{裕度}、DCM 相对参考的\emph{误差}、角动量的\emph{峰值}——不仅是经典 WBC 的监控面板，也是后续强化学习运动控制（P8）设计平衡类奖励时\emph{现成}的惩罚项：惩罚 ZMP 裕度过小、惩罚 DCM 误差过大、惩罚角动量峰值过高，就把“别摔、别甩过头”翻译成了可微的奖励。学习方法之所以大量借鉴本章概念，正因为这些量已经把“平衡好不好”量化成了带物理单位的数。详见 P8 的奖励设计章。
\end{note}

\begin{pitfall}[误区：CMM 在质心系，世界系角动量参考要先变换]\label{pit:hmw-cmm-frame}
Pinocchio \texttt{ccrba} 返回的 $\mathbf A_G$ 是在\emph{质心坐标系}中表达的。若你在\emph{世界系}里定义角动量参考 $\dot{\mathbf k}_G^d$（例如“让躯干绕世界竖直轴的角动量归零”），必须先把它变换到质心系再代入 \cref{eq:hmw-angular-task}，否则任务方程左右不在同一坐标系、QP 会朝错误方向使劲。注意：两系的\emph{平移}部分不影响角动量（角动量已是关于质心算的、与参考点平移无关），但\emph{旋转}部分必须一致。这与 \cref{ch:legged-biped} 机身角任务“旋转误差要在世界系算”（\cref{eq:lb-base-angular} 的 \texttt{rotationErrorInWorld}）是同一类“切空间/坐标系错配”陷阱。
\end{pitfall}

\begin{practice}
\begin{enumerate}
  \item 由 CMM 定义 \cref{eq:umm-cmm-def} 出发，独立推出角动量任务 \cref{eq:hmw-angular-task}，写清 Jacobian 与右端 $\mathbf b$ 各是什么；再代入“零角动量”参考 $\dot{\mathbf k}_G^d=-\mathbf K_k\mathbf A_G^{\mathrm{ang}}\mathbf v$，化简出含 $\mathbf K_k$ 的最终任务式。
  \item （\cref{ins:hmw-no-momentum-task} 的反事实）设计一个思想实验：人形单支撑站立，侧向受 $20\,\mathrm{N}$ 阶跃推力。分别在“有/无角动量任务”下，定性画出躯干倾角、$\mathbf k_G$、ZMP 三条随时间的曲线，标出无角动量任务时 ZMP 何时出界。
  \item （跨章综合）结合 \cref{ch:force-wbc} 的 WBC 框架与本节的角动量任务，回答：四足 WBC（如 MIT Cheetah 风格）与人形动量 WBC（IHMC 风格）在 QP \emph{目标函数}上的核心区别是什么？为什么四足通常\emph{不需要}角动量任务、人形\emph{必须}要？（提示：质量分布、上体高度、臂长、单步冲击。）
\end{enumerate}
\end{practice}

% ════════════════════════════════════════════════════════════════
\section{双支撑过渡：ZMP 转移、力 ramp 与等效 \texorpdfstring{$\omega_0$}{omega0}}
\label{sec:hmw-ds}

\paragraph{动机：被“瞬时切换”假设掩盖的那一瞬间。}
\cref{ch:legged-biped} 的解析落脚律把脚印切换当成\emph{瞬时}的——支撑脚从左“啪”地变右。但真实人形两脚有一段同时着地的\textbf{双支撑相}（double support, DS），典型占步态周期 $20\%$–$30\%$。这一段不是可有可无的过渡，它恰恰是“过渡时刻力矩跳变”这一常见故障的发生地，也是单支撑解析公式失准的根源。本节讲清它为何必须平滑、怎么平滑。

\subsection{ZMP 参考的连续转移}
\label{subsec:hmw-ds-zmp}

\paragraph{问题出在哪：阶跃 ZMP $\Rightarrow$ 阶跃力矩。}
单支撑相里 ZMP 只能在当前支撑脚的脚掌内；双支撑相里支撑多边形是\emph{两脚的外包络}，ZMP 可在两脚之间。如果规划在相位切换瞬间让 ZMP 参考从前脚中心\emph{直接跳}到后脚中心，那么由 LIP 关系 $\ddot x=\omega_0^2(x-p_{\mathrm{zmp}})$（\cref{eq:lb-lip}），$p_{\mathrm{zmp}}$ 阶跃 $\Rightarrow$ 质心加速度 $\ddot x$ 阶跃 $\Rightarrow$ 接触力阶跃 $\Rightarrow$ 关节力矩跳变。实机表现为过渡时刻“咯噔”一下，重则激起结构振动。

\paragraph{解法：让 ZMP 参考在 DS 相内连续转移。}
让 ZMP 参考在双支撑相内\emph{连续地}从前支撑脚移到后支撑脚，而非阶跃。用一个从 $0$ 平滑增到 $1$ 的过渡相位标量 $s$ 做插值：

\begin{equation}
  \boxed{\;\mathbf p_{\mathrm{zmp}}^{\mathrm{ref}}(t)=(1-s(t))\,\mathbf p_{\mathrm{foot}}^{\mathrm{prev}}+s(t)\,\mathbf p_{\mathrm{foot}}^{\mathrm{next}},\qquad s:0\to 1\ \text{在 DS 相内}.\;}
  \label{eq:hmw-ds-zmp}
\end{equation}

若 $s(t)$ 只取\emph{线性}（$s=t/T_{\mathrm{ds}}$），则 ZMP 速度在 DS 两端会有折角（$\dot s$ 不连续）；改用\emph{三次或五次}多项式让 $s$ 在两端满足 $\dot s=0$（乃至五次再加 $\ddot s=0$），就能保证 ZMP 的速度（乃至加速度）连续，从而经 LIP 关系传出去的力矩也连续。\cref{fig:hmw-ds} 对比阶跃与平滑两种 ZMP 转移及其力矩后果。

\begin{figure}[ht]
\centering
\begin{tikzpicture}[scale=1.0,>=Stealth]
  % left: step transition
  \begin{scope}[shift={(0,0)}]
    \draw[->] (0,0)--(3.2,0) node[right,font=\scriptsize]{$t$};
    \draw[->] (0,-0.2)--(0,1.9) node[above,font=\scriptsize]{$p_{\mathrm{zmp}}^{\mathrm{ref}}$};
    % step
    \draw[main!70!black,very thick] (0.2,0.3)--(1.5,0.3);
    \draw[main!70!black,very thick,dashed] (1.5,0.3)--(1.5,1.5);
    \draw[main!70!black,very thick] (1.5,1.5)--(3.0,1.5);
    \node[font=\scriptsize,gray!55!black] at (0.85,0.05) {前脚};
    \node[font=\scriptsize,gray!55!black] at (2.4,1.25) {后脚};
    % torque spike
    \draw[third,-{Stealth[length=2mm]},thick] (1.5,1.7)--(1.5,2.15);
    \node[third,font=\scriptsize,align=center] at (1.5,2.45) {力矩\\跳变};
    \node[main!60!black,font=\small] at (1.5,-0.5) {阶跃转移（错）};
  \end{scope}
  % right: smooth transition
  \begin{scope}[shift={(5.0,0)}]
    \draw[->] (0,0)--(3.2,0) node[right,font=\scriptsize]{$t$};
    \draw[->] (0,-0.2)--(0,1.9) node[above,font=\scriptsize]{$p_{\mathrm{zmp}}^{\mathrm{ref}}$};
    \draw[second!60!black,very thick] (0.2,0.3)--(1.0,0.3);
    % smooth S-curve (cubic-ish) drawn with bezier
    \draw[second!60!black,very thick] (1.0,0.3) .. controls (1.5,0.3) and (1.5,1.5) .. (2.0,1.5);
    \draw[second!60!black,very thick] (2.0,1.5)--(3.0,1.5);
    \node[font=\scriptsize,gray!55!black] at (0.6,0.05) {前脚};
    \node[font=\scriptsize,gray!55!black] at (2.6,1.25) {后脚};
    % DS window
    \draw[gray!50,dashed] (1.0,0)--(1.0,1.7);
    \draw[gray!50,dashed] (2.0,0)--(2.0,1.7);
    \draw[gray!50,<->] (1.0,1.75)--(2.0,1.75) node[midway,above,font=\scriptsize,gray!55!black]{DS 相};
    \node[second!50!black,font=\small] at (1.5,-0.5) {三/五次平滑（对）};
  \end{scope}
\end{tikzpicture}
\caption{双支撑期 ZMP 参考转移。左：阶跃转移——ZMP 在相位切换瞬间从前脚跳到后脚，经 $\ddot x=\omega_0^2(x-p_{\mathrm{zmp}})$ 引出质心加速度、接触力、关节力矩的同步跳变。右：在 DS 相窗口内用三/五次多项式让相位标量 $s:0\to1$ 平滑过渡（两端 $\dot s=0$），ZMP 速度乃至加速度连续，力矩随之连续。代价是行走略慢；DS 相越长越平滑越鲁棒、越短越接近瞬时切换的理想假设。}
\label{fig:hmw-ds}
\end{figure}

\subsection{接触力的渐入渐出}
\label{subsec:hmw-ds-force}

与 ZMP 转移配套，两脚的法向力也要平滑交接，而非瞬间换载。前支撑脚的法向力 $f_n^{\mathrm{prev}}$ 从满载\emph{渐降到零}（ramp-out），后支撑脚 $f_n^{\mathrm{next}}$ 从零\emph{渐升到满载}（ramp-in），二者之和始终约等于 $mg$（要支住整机重量）。WBC 层通过给两个接触任务设\emph{时变权重}或\emph{时变力参考}实现这一点——这正是摆脚 touchdown 时“接触力渐入”在双支撑相的体现。

\begin{pitfall}[误区：把双支撑当成“两脚硬撑、随便分力”]\label{pit:hmw-ds-force-free}
双支撑相\emph{不是}“两脚都硬支撑、力可任意分配”。双支撑期间两脚接触力\emph{各自}受自己的摩擦锥、CoP 边界约束，且\emph{总和}必须实现质心动力学（合力托住整机、合力矩满足平衡）。把它当成无约束力分配，QP 会算出某只脚法向力为\emph{负}（拉地面）的非物理解——在仿真里若接触模型软则表现为穿透/弹跳，在实机上直接是打滑或翻边。\cref{ch:force-wbc} 的法向单边约束 $f_n\ge 0$ 在双支撑期\emph{两只脚都要加}，缺一不可。
\end{pitfall}

\subsection{为何单支撑解析 DCM 在双支撑期有 \texorpdfstring{$2$--$5\,\mathrm{cm}$}{2-5cm} 误差}
\label{subsec:hmw-ds-omega}

\cref{ch:legged-biped} 用单支撑的解析 DCM/捕获点公式（\cref{eq:lb-capture}、\cref{sec:lb-alip}）算落脚，这套公式默认 $\omega_0=\sqrt{g/z_c}$ 在整步内\emph{恒定}、且支撑点瞬时切换。把它直接用到双支撑期，会有约 $2$–$5\,\mathrm{cm}$ 的系统落脚误差。\footnote{此 $2$–$5\,\mathrm{cm}$ 是\emph{经验量级}而非普适常数：它随机器人尺度、DS 相时长、行走速度而变（已联网核实——DCM/ZMP 文献普遍把瞬时切换假设带来的 DS 期偏差当作随步时/落脚/相时长变化的配置相关量，量级在厘米级，如初始 DCM 偏移 $>0.01\,\mathrm m$ 即足以在过渡期引出显著 ZMP 误差，参见 Englsberger 等连续 DS 轨迹生成的动机\cite{englsberger2015dcm}），切勿当作精确指标。其物理根（移动支撑点 $\Rightarrow$ 时变等效 $\omega_0$）见 \cref{ins:hmw-ds-omega}。}

\begin{insight}[误差的根：双支撑期“等效支撑点”在两脚间移动 $\Rightarrow$ 等效 $\omega_0$ 变了]\label{ins:hmw-ds-omega}
误差的物理根源藏在 $\omega_0$ 里。单支撑期支撑点固定，$\omega_0=\sqrt{g/z_c}$ 是个定值，DCM 动力学 $\dot{\boldsymbol\xi}=\omega_0(\boldsymbol\xi-\mathbf p)$ 干净成立。但双支撑期间，等效支撑点（ZMP 参考）正按 \cref{eq:hmw-ds-zmp} 在两脚之间\emph{移动}，而支撑力也在两脚间转移——这等价于一个\emph{时变}的有效倒立摆，其等效 $\omega_0$ 与单支撑值\emph{不同}。用单支撑的恒定 $\omega_0$ 去外推（指数因子 $e^{\omega_0 T}$）就会系统性偏差，落脚偏出几厘米。工程上两种务实处理：(1) 把 DS 期当作 ZMP 在两脚间移动的\emph{已知参考}，用分段 DCM 反向积分算 DCM 参考；(2) 简化系统直接缩短 DS 相、用单支撑公式近似——代价就是上面那 $2$–$5\,\mathrm{cm}$ 误差，对慢速行走可接受。\textbf{DS 相时长是一个直接体现“鲁棒性 vs 动态性”的旋钮：长则过渡平滑、力矩干净、但行走慢、不像人；短则接近瞬时切换理想、但对力矩平滑与状态估计要求更高。}
\end{insight}

\begin{practice}
\begin{enumerate}
  \item 写出一个满足 $s(0)=0,s(T_{\mathrm{ds}})=1,\dot s(0)=\dot s(T_{\mathrm{ds}})=0$ 的三次多项式 $s(t)$；再加上 $\ddot s(0)=\ddot s(T_{\mathrm{ds}})=0$ 给出五次多项式。把它代入 \cref{eq:hmw-ds-zmp}，说明为什么五次比三次更能压低力矩过渡的“抖”。
  \item （\cref{pit:hmw-ds-force-free}）设 DS 相内 $f_n^{\mathrm{prev}}+f_n^{\mathrm{next}}=mg$ 恒成立、且两者都用线性 ramp。写出某一时刻若把全部 $mg$ 仍压在前脚（ramp 未生效）会导致后脚 $f_n^{\mathrm{next}}$ 取什么值——它为何非物理？
  \item （概念，承 \cref{ins:hmw-ds-omega}）若把 DS 相时长缩到接近零，本章这套“连续 ZMP 转移 + 力 ramp”是否还必要？这把人形行走推向 \cref{ch:legged-biped} 哪个理想假设？代价是什么？
\end{enumerate}
\end{practice}

% ════════════════════════════════════════════════════════════════
\section*{本章小结}
\addcontentsline{toc}{section}{本章小结}

本章在 \cref{ch:legged-biped}（点足落脚平衡）、\cref{ch:unified-multimodal-mpc}（CMM/角动量守恒）、\cref{ch:force-wbc}/\cref{ch:quad_wbc}（WBC-QP）之上，补齐了人形\emph{特有}的两件事——\textbf{脚掌带来踝/髋臂两档余量、上体可主动用角动量}——并把它们落进全身 QP。一句话主线：\textbf{脚掌让“平衡只能迈步”放宽为“踝→髋臂→迈步”三级，CMP 把这三级几何化、飞轮把“甩臂买时间”定量化、CMM 角行块把它落成 QP 里的角动量任务；而两脚交接的双支撑相要靠连续 ZMP 转移与力 ramp 才不抖。}

\begin{itemize}
  \item \textbf{CMP 与三级仪表盘}（\cref{sec:hmw-cmp}）：$\mathbf r_{\mathrm{CMP}}=\mathbf r_{\mathrm{ZMP}}+\frac{1}{f_z}[\tau_{G,y};-\tau_{G,x}]$ 把角动量效应折算成长度；ZMP 必在脚内、CMP 可出脚、捕获点指迈步，三点各答一问、合成平衡仪表盘；CMP 出脚有代价（角动量饱和）。
  \item \textbf{飞轮等效捕获点}（\cref{sec:hmw-flywheel}）：$\boldsymbol\xi_{eq}=\boldsymbol\xi-\Delta\mathbf k_G/(m\omega_0 z_c)$，角动量把等效捕获点拉近、暂时撑大“不迈步可恢复区”；守恒与摆臂救场分属腾空/接触两相，靠接触力矩这道闸门分开（臂作反作用轮，去重指向 \cref{sec:umm-momentum}）。
  \item \textbf{角动量任务与人形任务栈}（\cref{sec:hmw-momentum-wbc}）：由 CMM 求导得 $\mathbf A_G^{\mathrm{ang}}\dot{\mathbf v}=\dot{\mathbf k}_G^d-\dot{\mathbf A}_G^{\mathrm{ang}}\mathbf v$；人形 TSID/HQP 栈 $=$ 四足栈 $+$ 这一条角动量任务（IHMC/TALOS 动量优先）；无此任务则侧推下身越来越斜直到 ZMP 出界；ZMP 裕度/DCM 误差/角动量峰值是天然奖励信号（连 P8）。
  \item \textbf{双支撑过渡}（\cref{sec:hmw-ds}）：$\mathbf p_{\mathrm{zmp}}=(1-s)\mathbf p_{\mathrm{foot}}^{\mathrm{prev}}+s\,\mathbf p_{\mathrm{foot}}^{\mathrm{next}}$ 用三/五次多项式保力矩连续；接触力 ramp-in/out 且和恒为 $mg$；单支撑解析 DCM 在 DS 期 $2$–$5\,\mathrm{cm}$ 误差的根是等效 $\omega_0$ 随移动支撑点变化。
\end{itemize}

\begin{table}[H]\centering\small
\caption{本章主要符号}
\begin{tabular}{lll}
\toprule
符号 & 含义 & 出处 \\
\midrule
$\mathbf r_{\mathrm{ZMP}}$ & 零力矩点（必在支撑域内） & \cref{ins:hmw-three-points} \\
$\mathbf r_{\mathrm{CMP}}$ & 质心力矩枢轴点（可出脚） & \cref{eq:hmw-cmp-def} \\
$\boldsymbol\tau_G=\dot{\mathbf k}_G$ & 绕质心力矩（$=$ 角动量变化率） & \cref{eq:hmw-cmp-def} \\
$f_z$ & 地面法向力 & \cref{eq:hmw-cmp-def} \\
$\boldsymbol\xi$ / $\boldsymbol\xi_{eq}$ & 捕获点 / 飞轮等效捕获点 & \cref{eq:lb-capture},\cref{eq:hmw-flywheel-xi} \\
$\Delta\mathbf k_G$ & 飞轮（臂/躯干）提供的角动量 & \cref{eq:hmw-flywheel-xi} \\
$\mathbf A_G$ / $\mathbf A_G^{\mathrm{ang}}$ & CMM / 其角行块（后三行） & \cref{eq:umm-cmm-def},\cref{eq:hmw-angular-task} \\
$\dot{\mathbf k}_G^d$ & 角动量任务期望变化率 & \cref{eq:hmw-angular-task} \\
$\omega_0=\sqrt{g/z_c}$ & LIP 自然频率（$=$ \cref{ch:legged-biped} 的 $\omega$） & \cref{eq:lb-lip} \\
$s:0\to1$ & 双支撑过渡相位标量 & \cref{eq:hmw-ds-zmp} \\
\bottomrule
\end{tabular}
\end{table}

\begin{table}[H]\centering\small
\caption{公式速查表}
\begin{tabular}{lll}
\toprule
公式 / 结论 & 一句话说明 & 对应 \\
\midrule
CMP 定义 \cref{eq:hmw-cmp-def} & $\mathbf r_{\mathrm{CMP}}=\mathbf r_{\mathrm{ZMP}}+\frac{1}{f_z}[\tau_{G,y};-\tau_{G,x}]$ & \cref{ins:hmw-three-points} \\
三级策略仪表盘 & 踝(CMP=ZMP)/髋臂(CMP出脚)/迈步(捕获点出脚) & \cref{tab:hmw-strategies} \\
飞轮等效捕获点 \cref{eq:hmw-flywheel-xi} & 角动量把 $\boldsymbol\xi$ 拉近、暂撑大可恢复区 & \cref{ins:hmw-flywheel-region} \\
角动量任务 \cref{eq:hmw-angular-task} & $\mathbf A_G^{\mathrm{ang}}\dot{\mathbf v}=\dot{\mathbf k}_G^d-\dot{\mathbf A}_G^{\mathrm{ang}}\mathbf v$ & \cref{ins:hmw-no-momentum-task} \\
人形任务栈 \cref{tab:hmw-task-stack} & $=$ 四足栈 $+$ 角动量任务 & \cref{ins:hmw-vs-quad} \\
DS 期 ZMP 转移 \cref{eq:hmw-ds-zmp} & 三/五次平滑保力矩连续 & \cref{fig:hmw-ds} \\
DS 期 DCM 误差 & $2$–$5\,\mathrm{cm}$，根在等效 $\omega_0$ 变化 & \cref{ins:hmw-ds-omega} \\
\bottomrule
\end{tabular}
\end{table}

\begin{table}[H]\centering\small
\caption{知识点总表（行星号为难度）}
\begin{tabular}{clll}
\toprule
\# & 知识点 & 核心要点 & 难度 \\
\midrule
1 & CMP 定义 & 角动量折算长度、可出脚 & \(\bigstar\bigstar\bigstar\) \\
2 & 三级仪表盘 & ZMP/CMP/捕获点各答一问 & \(\bigstar\bigstar\bigstar\) \\
3 & 飞轮等效捕获点 & 角动量买时间、撑大可恢复区 & \(\bigstar\bigstar\bigstar\) \\
4 & 守恒不矛盾 & 接触力矩是角动量进出口 & \(\bigstar\bigstar\bigstar\) \\
5 & 角动量任务 & 由 CMM 角行求导而来 & \(\bigstar\bigstar\bigstar\bigstar\) \\
6 & 人形 vs 四足栈 & 多一条角动量任务 & \(\bigstar\bigstar\bigstar\) \\
7 & 双支撑 ZMP 转移 & 三/五次保力矩连续 & \(\bigstar\bigstar\bigstar\) \\
8 & DS 期等效 $\omega_0$ & 移动支撑点 $\Rightarrow$ DCM 误差 & \(\bigstar\bigstar\bigstar\bigstar\) \\
\bottomrule
\end{tabular}
\end{table}

% ════════════════════════════════════════════════════════════════
\section*{延伸阅读}
\begin{itemize}
  \item \textbf{CMP 与恢复策略}：质心力矩枢轴点与踝/髋策略的几何由 Popovic、Goswami、Herr 等系统化\cite{popovic2005cmp}；捕获点与飞轮模型的奠基见 Pratt 等\cite{pratt2006capture}（本章 \cref{eq:hmw-flywheel-xi} 同脉络）；可捕获性的分级理论见 Koolen 等\cite{koolen2012capturability}。
  \item \textbf{动量 WBC 架构}：IHMC Atlas 的动量优先框架见 Koolen、Bertrand 等\cite{koolen2016atlas}（本章 \cref{ins:hmw-no-momentum-task}/\cref{ins:hmw-vs-quad} 的来源）；无关节力矩传感器的力矩级控制（在 HRP-2 上实现、TSID 谱系）见 Del Prete/Mansard 等\cite{delprete2016tsid}；四足侧的脉冲控制 WBIC 对照见 Kim 等\cite{kim2019wbic}。
  \item \textbf{DCM/VRP 与连续行走}：DCM 的三维系统化与 VRP 见 Englsberger 等\cite{englsberger2015dcm}（亦见 \cref{ch:legged-biped} 延伸阅读）。
  \item \textbf{质心动量与 CMM}：质心动量矩阵/CCRBA 的奠基见 Orin 等\cite{orin2013centroidal}；本书 CMM 详述见 \cref{ch:unified-multimodal-mpc}。
\end{itemize}

\section*{本章与后续章节的关系}
\begin{table}[H]\centering\small
\begin{tabular}{lll}
\toprule
相关章节 & 关系 & 本章接口 \\
\midrule
\cref{ch:legged-biped} 欠驱点足 & 提供 LIP/DCM/捕获点地基，本章放宽“只能迈步” & \cref{eq:lb-capture},\cref{ins:hmw-relax-stepping} \\
\cref{ch:unified-multimodal-mpc} 统一动力学 & 提供 CMM 与角动量守恒，本章组装角动量任务 & \cref{eq:umm-cmm-def},\cref{eq:hmw-angular-task} \\
\cref{ch:force-wbc} 力控与 WBC & 提供 WBC-QP 框架，本章任务插入其中 & \cref{eq:fw-wbcqp},\cref{eq:hmw-torque-recover} \\
\cref{ch:quad_wbc} 四足 WBC & 任务栈对照，人形多一条角动量任务 & \cref{eq:qw-wqp},\cref{tab:hmw-task-stack} \\
\cref{ch:lie} 李群与李代数 & 角动量参考的坐标系变换 & \cref{pit:hmw-cmm-frame} \\
\bottomrule
\end{tabular}
\end{table}

\section*{故障排查手册}
\begin{enumerate}
  \item \textbf{症状}：人形受侧推后\emph{躯干越来越斜}、最后单支撑时 ZMP 出界翻倒，力矩跟踪却“很准”。\textbf{先查}：任务栈里\emph{有没有}角动量任务（\cref{ins:hmw-no-momentum-task}）；若无，角动量在上体无人约束、必持续累积——补上 \cref{eq:hmw-angular-task} 这条任务并放高优先级。
  \item \textbf{症状}：手臂去抓远处重物时，\emph{身体平衡变差}、ZMP 逼近边界。\textbf{原因}：把臂末端任务与 CoM/角动量\emph{同权加权}，臂的大 Jacobian 引发优先级泄漏（\cref{ins:hmw-hqp-momentum}）。\textbf{对策}：改用 HQP 把 CoM+角动量放高层、操作任务进零空间（\cref{sec:qw-nullspace}）；先稳后抓、再逐步提操作权重。
  \item \textbf{症状}：角动量任务一加，\emph{躯干/臂高频抖动}。\textbf{原因}：角动量增益 $\mathbf K_k$ 过大，快速压平角动量需大关节加速度（\cref{subsec:hmw-momentum-task}）。\textbf{对策}：减小 $\mathbf K_k$；确认参考 $\dot{\mathbf k}_G^d$ 已变换到质心系（\cref{pit:hmw-cmm-frame}）。
  \item \textbf{症状}：双支撑相位切换瞬间力矩\emph{“咯噔”跳变}、激起结构振动。\textbf{原因}：ZMP 参考阶跃转移（\cref{subsec:hmw-ds-zmp}）。\textbf{对策}：用 \cref{eq:hmw-ds-zmp} 的三/五次平滑 $s(t)$；配套两脚力 ramp-in/out（\cref{subsec:hmw-ds-force}）。
  \item \textbf{症状}：双支撑期 QP 算出某只脚法向力为\emph{负}（拉地面）。\textbf{原因}：把 DS 当成无约束力分配、漏了对\emph{两只}脚都加 $f_n\ge 0$（\cref{pit:hmw-ds-force-free}）；也可能叠加身体已失稳。\textbf{对策}：两脚都加法向单边约束，并检查 CoP 是否各自在脚内。
  \item \textbf{症状}：用单支撑解析 DCM 算落脚，慢走时\emph{系统性偏出几厘米}。\textbf{原因}：双支撑期等效 $\omega_0$ 随移动支撑点变化、与单支撑值不同（\cref{ins:hmw-ds-omega}）。\textbf{对策}：要么缩短 DS 相接受 $2$–$5\,\mathrm{cm}$ 误差，要么按 ZMP 在两脚间的已知参考做分段 DCM 反向积分。
  \item \textbf{症状}：实机 CoP（传感器）与控制器算的 ZMP \emph{对不上}。\textbf{原因}：二者身份不同，倾斜地/足底扭矩下本就可能不同（\cref{pit:hmw-zmp-cop}）。\textbf{对策}：分别记录两者，其差是接触建模误差的体检指标，不要当成 bug 强行对齐。
\end{enumerate}
